\section{Discussion}
\label{sec:discussion}

The practical reading of the wall is that computed values are an unguarded trust boundary in deployed systems. Agent pipelines pass computed results from step to step \citep{yao2023react,kapoor2024agents}, and chain-of-thought monitors watch for signs of doubt \citep{baker2025monitoring,korbak2025monitorability}. Our measurements say that a corrupted computed value crosses both defenses silently: no model rechecks it and almost none remarks on it. Nor is the trust self-trust that accurate provenance labels might break. Absorption of computed errors is unchanged when the trace is attributed to another model or to a human (Appendix~\ref{app:attribution}), so a wrong value entering from a tool call or a collaborating agent is carried exactly as the model's own would be. The failure mode this predicts is specific: not noisy degradation but confident, arithmetically consistent propagation of a single wrong intermediate, invisible to any monitor that reads the model's own commentary. It is the within-trace analog of hallucinations snowballing into confident elaboration \citep{zhang2023snowball}.

We interpret the default as a learned prior rather than a capability gap, and mark this paragraph as interpretation. In ordinary text, a value written as the result of a computation is almost always correct, so trusting it is the right prediction on the training distribution. Text rarely shows an author stopping mid-derivation to re-verify a line \citep{mccoy2024embers,prystawski2023why}. The improvements that separate our strongest and weakest models came from training that rewards catching visible inconsistencies, and the readable side of our measurements improved from near zero to near perfect accordingly. Nothing in that training touches the computed side, and the computed side did not move. A prior installed by the distribution and untouched by the objective would look exactly like the flat ceiling we observe, including its indifference to instruction. The instruction experiment shows the model reads the command to recompute and continues predicting text in which computed values are settled.

We also looked for the missing check inside one open model, and found it present but unused. On Qwen2.5-7B we trained linear probes on minimal pairs of traces that differ only in the planted value. We read the residual stream at the first position after the planted line, where the tokens are identical in both members of the pair. The probe is at chance through 25 of the model's 28 layers. It separates wrong from right values only in the last three (held-out AUC 0.80, 0.77, and 0.98 at layers 26, 27, and 28), and at the planted value itself it reaches 1.00 at layer 27. The model computes that the committed value is wrong, but late. Steering along these directions did not turn the computation into a check. We tried 150 conditions: single layers, layer bands, and the late layers where the signal lives, both signs, and doses up to the direction's natural magnitude. In every condition that left the model able to continue unperturbed traces, absorption stayed between 0.78 and 0.95 against a 0.90 baseline, and explicit checking never exceeded 5 percent of continuations. The only doses that produced more checking language had also destroyed ordinary continuation. We read this as the representational side of deference: the information exists, in a form the generator does not consult, and amplifying it does not make the generator consult it. A fuller account of this representation is future work.

The same reading says what should fix it, and our protocol supplies the test. Dedicated step verifiers already learn to catch computation errors \citep{cobbe2021verifiers,uesato2022process,lightman2023,wang2024mathshepherd,mcaleese2024critics}, so the checking behavior is learnable; what remains unmeasured is whether training installs the generating model's propensity to deploy it mid-trace. Training against planted errors is the untested lever on this task, though not untried elsewhere. Reinforcement learning on corrupted-prefix samples improves outcome-level recovery \citep{rlflawed2025}, while supervised fine-tuning on correction traces can fail outright, with corrected models parroting the original mistake \citep{sfterr2025}. Neither has been evaluated against unprompted verification of committed values under an audited instrument, which is the experiment our protocol makes possible. System-side re-execution of computable content is the immediate engineering workaround, since for programs and arithmetic the check is exact and cheap, with sampling-based consistency checks \citep{wang2023selfconsistency,manakul2023selfcheckgpt} the fallback where it is not. Either way, the planted-error protocol gives a mechanical pass-fail: a model breaches the wall when its absorption of computed errors leaves the ceiling, a number any lab can measure in an afternoon.

Four limitations bound the claims. Our task is straight-line integer programs, chosen because ground truth is mechanical; computed content in freer forms, such as word problems or code with control flow, is untested. Reasoning-channel models are tested in Section~\ref{sec:reasoning} under their own deliberation interfaces rather than the roster protocol; one further frontier model could not be measured at all because its safety classifier refuses the task format. Five of our thirteen models are measured through the user-turn method rather than prefilling, and one through a legacy completion endpoint, calibrated but not identical. And our claim is behavioral: probing work suggests models internally represent step correctness they do not act on, so the wall describes what models do, not what they encode.
